\documentclass[12pt,a4paper]{article}

\usepackage[utf8]{inputenc}
\usepackage[T1]{fontenc}
\usepackage[margin=2.5cm]{geometry}
\usepackage{hyperref}
\usepackage{microtype}
\usepackage{parskip}
\usepackage{lmodern}

\hypersetup{colorlinks=true, urlcolor=blue!60!black}

\pagestyle{empty}

\begin{document}

\begin{flushright}
Kundai Farai Sachikonye \\
Auf der Lichtung 19, 82178 Puchheim, Germany \\
\href{https://github.com/fullscreen-triangle}{github.com/fullscreen-triangle} \\
\texttt{kundai.sachikonye@bitspark.com} \\
(+49) 1626386344 \\[6pt]
1 June 2026
\end{flushright}

\vspace{1em}

Deutsches Zentrum f\"ur Astrophysik (DZA) \\
Department Algorithmik \\
G\"orlitz, Germany \\
Job-ID: D03-26-04

\vspace{1.5em}

\textbf{Re: Research Scientist --- Real-Time Data Processing \& Algorithm
Development (D03-26-04)}

\vspace{1em}

Dear Hiring Committee,

I am applying for the Research Scientist position in real-time data processing
and algorithm development. The core of the role --- designing algorithms for
real-time data extraction and building scalable, high-performance pipelines for
extremely large data streams --- is the core of what I have spent the last
several years doing. My domain has been computational biology rather than
astrophysics, but the engineering problem is the same: extract signal from data
streams too large to hold or process by conventional means, in time, at scale.

\textbf{High-performance pipelines for very large data streams.}
I build and deploy these systems. \href{https://github.com/fullscreen-triangle/lavoisier}{Lavoisier}
is a distributed analysis pipeline that processes mass-spectrometry data streams
exceeding 100\,GB: memory-mapped I/O so dataset size is bounded by storage rather
than memory, Ray-distributed execution across HPC nodes, a structured HDF5 output
schema, and a dual numerical/CNN path that reaches 98.9\,\% feature-extraction
accuracy against NIST references. \href{https://github.com/fullscreen-triangle/gospel}{Gospel}
sustains $10^6$ records per second of streaming throughput with Bayesian-network
and mixture-model classification, validated against 1000~Genomes, gnomAD, and
UK~Biobank. \href{https://github.com/fullscreen-triangle/four-sided-triangle}{Four-Sided-Triangle}
is a multi-stage engine with a Rust core (10--50$\times$ over the reference
implementation) built for parallel, horizontally-scaled execution. These are the
methods the position calls for, demonstrated on production data.

\textbf{The required expertise areas.}
The posting asks for demonstrated expertise in at least one of several areas; my
work spans most of them. Large-scale data analytics and real-time stream
processing are the substance of the systems above. Machine learning and AI run
through all of them (PyTorch CNNs and sequence models, Bayesian inference,
GPU-accelerated computation). Parallel and high-performance computing --- Ray,
Dask, SLURM-compatible design, Docker/Kubernetes, memory-mapped streaming --- is
the infrastructure I work in daily. Python is my primary language, with Rust for
performance-critical cores.

\textbf{Astrophysical signal detection.}
The detection, classification, and characterisation of signals in large data
streams is a problem I have been drawn to independently, including work applying
these stream-processing and signal-extraction methods to astronomical data ---
publicly available alongside my other systems. I would welcome bringing the
real-time and HPC methods I have built in one data-intensive domain to the
astrophysical data streams the DZA works with, and contributing to the
algorithmic foundations of that processing.

\textbf{Background.}
I hold an MSc in Computational Biology (Universit\"at T\"ubingen) and an MSc in
Pharmaceutical Biotechnology (HAW Hamburg), and I conducted doctoral research in
Computational Lipidomics at TUM (2019--2021) before transitioning to independent
research, where my output has been validated, publicly documented data-processing
systems. I work independently in international, English-language settings; English
is native and my German is conversational. I am resident in Germany with full
work authorisation and available at the earliest possible date.

I would welcome the opportunity to discuss how this background in real-time,
large-scale data processing could contribute to the DZA's algorithmics work.

\vspace{1.5em}
Yours sincerely,

\vspace{2.5em}
\textbf{Kundai Farai Sachikonye}

\end{document}
